\documentclass[11pt]{amsart}
\usepackage{calc,amssymb,amsmath,ulem,stmaryrd,fullpage}
\usepackage{enumerate}
\usepackage{alltt}
\RequirePackage[dvipsnames,usenames]{color}
\let\mffont=\sf
\normalem
%\input{kmacros3.sty}
\input{xy}
\xyoption{all}
\usepackage{hyperref}
\usepackage{graphicx}
\usepackage[all,cmtip]{xy}
\usepackage{verbatim}
\usepackage[left=0.95in,top=0.95in,right=0.95in,bottom=0.95in]{geometry}
\newcommand{\tauCohomology}{T}
\newcommand{\FCohomology}{S}


%\theoremstyle{theorem}
%\newtheorem*{mainthm*}{Main Theorem}

\newcommand{\note}[1]{\marginpar{\sffamily\tiny #1}}

\def\todo#1{\textcolor{Mahogany}%
{\sffamily\footnotesize\newline{\color{Mahogany}\fbox{\parbox{\textwidth-15pt}{#1}}}\newline}}

\renewcommand{\O}{\mathcal O}
\newcommand{\ie}{{\itshape i.e.} }
\newcommand{\roundup}[1]{\lceil #1 \rceil}
\newcommand{\rounddown}[1]{\lfloor #1 \rfloor}
\renewcommand{\to}[1][]{\xrightarrow{\ #1\ }}
\newcommand{\onto}[1][]{\protect{\xrightarrow{\ #1\ }\hspace{-0.8em}\rightarrow}}
\newcommand{\into}[1][]{\lhook \joinrel \xrightarrow{\ #1\ }}
%\newcommand{DBDual}[1]{{\underline{\omega}_{#1}}}
\newcommand{\DBDual}[1]{{{\uuline \omega}{}^{\mydot}_{#1}}}
\newcommand{\Tt}{{\mathfrak{T}}}
%\newcommand{\bn}{{\mathfrak{n}} }
\newcommand{\sol}[1]{ \vskip 6pt{\color{blue} {\bf Solution:} #1} \vskip 6pt}

\newcommand{\bR}{{\mathbb{R}}}
\newcommand{\va}{{\vec{a}}}
\newcommand{\vb}{{\vec{b}}}
\newcommand{\vzero}{{\vec{0}}}
\newcommand{\vi}{{\vec{i}}}
\newcommand{\vj}{{\vec{j}}}
\newcommand{\vk}{{\vec{k}}}
\newcommand{\vh}{{\vec{h}}}
\newcommand{\vr}{{\vec{r}}}
\newcommand{\vu}{{\vec{u}}}
\newcommand{\vv}{{\vec{v}}}
\newcommand{\vw}{{\vec{w}}}
\newcommand{\vx}{{\vec{x}}}
\newcommand{\vy}{{\vec{y}}}
\newcommand{\vz}{{\vec{z}}}
\newcommand{\vT}{{\vec{T}}}
\newcommand{\vN}{{\vec{N}}}
\newcommand{\vF}{{\vec{F}}}
\newcommand{\vG}{{\vec{G}}}
\newcommand{\bF}{\mathbb{F}}
\newcommand{\bQ}{\mathbb{Q}}
\newcommand{\bZ}{\mathbb{Z}}
\newcommand{\bC}{\mathbb{C}}
\newcommand{\Inn}{\mathrm{Inn}}
\newcommand{\Aut}{\mathrm{Aut}}
\newcommand{\Hom}{\mathrm{Hom}}
\newcommand{\End}{\mathrm{End}}

\begin{document}
\title{HW \#5 -- Math 6310\\Fall 2019}
\author{Due: Wednesday, November 13th}
\maketitle
\noindent
{\bf 1.}  Suppose $D$ is a PID and $M$ is a finitely generated $D$-module.  Prove that 
\[
    M \cong D/\langle f_1(x) \rangle \oplus \dots \oplus D/\langle f_l(x) \rangle
\]
where $f_i | f_{i-1}$ for $i = 2, \dots, l$ and where the $f_i$ are monic polynomials.  Furthermore, show that this decomposition is unique.
\vfill
In many texts, rational canonical form also means that associated polynomials of the block diagonal divide each other in the way above (or maybe $f_i | f_{i+1}$).
\newpage
\noindent
{\bf 2.}  Determine the possible characteristic and minimal polynomials of an $n\times n$ matrix over $\bC$ that has rank $1$.
\newpage
\noindent
{\bf 3.}  Let $G$ be a finite Abelian group under addition and let $n > 0$ be an integer.  Show that the map $\phi : G \to G$ which sends $x \mapsto n x$ is a homomorphism, further show it is an isomorphism if and only if $\gcd(|G|, n) = 1$.
\newpage
\noindent
{\bf 4.}  Determine, up to conjugacy, all $3\times 3$ matrices $M$ over $\bQ$ that satisfy $M^3=2M^2$.
\newpage
\noindent
{\bf 5.}  Suppose that $k$ is the field with 3 elements and let $R = k[x]$.  Identify eleven different (up to isomorphism) $R$-modules $M$ such that $|M| = 9$. 
\newpage
\noindent
{\bf 6.}  Determine up to similarity, all \emph{real} $5 \times 5$-matrices with characteristic polynomial $x(x^2+1)^2$.  
\newpage
\noindent
{\bf 7.}  Let $M$ be a $3 \times 3$ complex matrix with $M^6=M^4$ and $M^4+M^2=2M^3$. Determine the possible Jordan forms of $M$.
\newpage
\noindent
{\bf 8.}  Write down the rational and Jordan canonical forms of the following matrix (viewed over $\bC$)
\[
    \left[
        \begin{array}{cccc}
            1 & 0 & 0 & 0\\
            0 & 1 & 0 & 0\\
            -2 & -2 & 0 & 1\\
            -2 & 0 & -1 & -2
        \end{array}
    \right].
\]
\emph{Hint:}  Let $V = \bC^4$ and view $\bC[x]$ acting on $V$ where $x$ acts as that matrix.  The 4 $\bC$-basis elements generate $V$, but there are relations between them when viewed over $\bC[x]$.
\newpage
\noindent
{\bf 9.}  Compute the characteristic polynomial, the minimal polynomial, the Jordan normal form, and the rational canonical form of the matrix.
\[
    \left[
        \begin{array}{ccc}
            3 & 1 & -1 \\
            2 & 2 & -1 \\
            2 & 2 & 0
        \end{array}
    \right].
\]


\end{document}

